\documentclass[10pt]{article}
\usepackage[T1]{fontenc}
\usepackage{lmodern,microtype,amsmath,amsthm,amssymb,booktabs,mathtools}
\usepackage[margin=1.05in]{geometry}
\usepackage[hidelinks]{hyperref}
\newtheorem{theorem}{Theorem}
\newtheorem{lemma}[theorem]{Lemma}
\newtheorem{corollary}[theorem]{Corollary}
\newcommand{\Z}{\mathbb Z}
\newcommand{\Q}{\mathbb Q}
\newcommand{\C}{\mathbb C}
\newcommand{\ord}{\operatorname{ord}}
\title{Periodic modes of Grover walks on finite abelian groups}
\author{}
\date{}
\begin{document}
\maketitle
\vspace{-2.3em}
\begin{abstract}
We classify the root-of-unity Fourier modes of the Grover walk with $d$
pairs of steps on an abelian group of odd prime-power exponent. When
$p\nmid d$, all periodic quadratic blocks come from characters taking
the same value, up to inversion, on all the steps. When $p\mid d$, exactly
two additional traces can occur. On $(\Z/N\Z)^d$, $N=p^h$, the
multiplicity of each primitive fourth root is a polynomial of degree $d/p$
in $N$, and that of each primitive sixth root has degree $\lfloor d/(2p)\rfloor$.
We give their leading coefficients and an exact example in dimension six.
For two pairs of steps we also give the complete classification at arbitrary
exponent. These results are unconditional.
\end{abstract}
\section{The Fourier blocks}
The proposition on periodic states on an odd rectangular grid in
the graph-tower section of \cite{repo} counts the periodic subspace of
the two-dimensional Grover walk with both side lengths odd. Its proof
reduces a relation among roots of unity modulo $2$. The question pursued
here is what changes when the number of step pairs and the prime in the
tower both vary. Akahori, Hayashi and Suzuki identify root-of-unity
occurrence with torsion zeros of a spectral resultant and suggest the
reduced periodic spectrum in Remark~6.5 \cite{ahs}. Our contribution is an
explicit character classification and a multiplicity law for these tori.
The two-pair test also follows from the rational-angle product equation
treated by Speyer \cite{speyer}, via the cosine addition formula. We include
its short valuation proof to cover arbitrary two-generator quotients.
The priority of the higher-dimensional statements has not been established
by an exhaustive literature search.

Let $G$ be a finite abelian group, and $s_1,\ldots,s_d$ a generating list,
$d\ge2$. Use one labelled pair of arcs for $s_i,-s_i$ even if their
endpoints coincide. For each character $\chi$ write $X_i=\chi(s_i)$ and
$C(x)=x+x^{-1}$. The Grover coin is $2uu^*-I$ on the $2d$ directions,
where $u=(1,\ldots,1)/\sqrt{2d}$. In the Fourier block the arc reversal
and shift combine to give the involution $S$ whose $i$th pair has matrix
$\left(\begin{smallmatrix}0&X_i\\X_i^{-1}&0\end{smallmatrix}\right)$.
The walk on this block is $U=S(2uu^*-I)$. Its characteristic polynomial is
\begin{equation}\label{eq:fourier}
 (z^2-1)^{d-1}(z^2-Tz+1),
 \qquad T=\frac1d\sum_{i=1}^d C(X_i).
\end{equation}
Indeed $\det(zI+S)=(z^2-1)^d$ and the rank-one determinant identity
multiplies this by $1-2u^*(zI+S)^{-1}Su$. Since
$(zI+S)^{-1}S=(zS-I)/(z^2-1)$ and $2u^*Su=T$, the identity follows
first away from $\pm1$, then polynomially everywhere.

The constant factor in the spectrum comprises $d-1$ copies of each of
$1,-1$ per character. We call the remaining degree-two factor the
\emph{quadratic block}. Both its eigenvalues have finite multiplicative
order together. Since the full walk is unitary, the periodic subspace
is the direct sum of all root-of-unity eigenspaces.

\begin{lemma}\label{l:int}
The quadratic block is periodic if and only if $T$ is an algebraic integer.
\end{lemma}
\begin{proof}
If $w$ is a root of unity, $w+w^{-1}$ is integral. Conversely,
the roots of $z^2-Tz+1$ are algebraic integers when $T$ is integral.
Every conjugate of $T$ is an average of numbers $C(Y_i)$ with
$|Y_i|=1$, hence is real and lies in $[-2,2]$. Every conjugate of either
root therefore has modulus one. Kronecker's theorem gives the claim:
the coefficients of the monic polynomials whose roots are the $n$th
powers of all conjugates are bounded integers, so these powers range over
a finite set and a nonzero root has finite order.
\end{proof}

\section{Odd prime-power exponent}
Let $p$ be an odd prime, $N=p^h$ with $h\ge1$, and suppose the exponent
of $G$ divides $N$. Write $H=\mu_p\subset\mu_N$, and let $c(x)$ be the
multiplicity of $x$ among the $2d$ entries $X_i,X_i^{-1}$. Call a tuple
\emph{aligned} if there is a $u\in\mu_N$ such that every $X_i$ belongs
to $\{u,u^{-1}\}$. For a coset $A$ let $\boldsymbol1_A$ denote its
indicator as a function on $\mu_N$.
\begin{theorem}\label{t:main}
A quadratic block is periodic exactly in the following cases.
\begin{enumerate}
\item The tuple is aligned. Its quadratic roots are $u,u^{-1}$, including
the double root $1$ at the trivial character.
\item For $k=0$ or $1$ there are nonnegative integers $n_A$ such that
\[
 c=kd\,\boldsymbol1_{\{1\}}+\sum_{A\in\mu_N/H} n_A\boldsymbol1_A .
\]
Then $p\mid d$ and $T=k$. The roots are the two primitive fourth roots for
$k=0$, and the two primitive sixth roots for $k=1$.
\end{enumerate}
Both additional traces occur on the standard $d$-torus of every side
length $p^h$ when $p\mid d$.
\end{theorem}
\begin{proof}
Let $s=N/p$ and $\zeta$ be a primitive $N$th root. Write
$c_j=c(\zeta^j)$ for $0\le j<N$. The ring of integers of $\Q(\zeta)$
is $\Z[\zeta]$ \cite[Chapter~2]{washington}. The cyclotomic relation
$\Phi_N(x)=1+x^s+\cdots+x^{(p-1)s}$ expresses the last power in
each list $\zeta^{a+js}$, $0\le j<p$, $0\le a<s$, as minus the
sum of its predecessors. The surviving powers form the integral power
basis. Thus the coordinates of $dT$ are
$c_{a+js}-c_{a+(p-1)s}$, $j<p-1$, and $T$ is integral precisely when
the counts are congruent modulo $d$ in each $H$-coset.

For $x\ne1$ at most one of $X_i,X_i^{-1}$ equals $x$, since $N$ is
odd. If $c(x)=d$ then all $X_i$ lie in $\{x,x^{-1}\}$.
For a nonaligned tuple we have $0\le c(x)<d$ away from $1$, so
congruence in each coset gives equality away from $1$. In $H$, let $r$
be the common value on its nonidentity elements.
Then $c(1)-r=kd$, where $0\le k\le2$ because $-d<c(1)-r\le2d$.
The case $k=2$ forces every $X_i=1$. The other two cases give the
asserted decomposition. Each coset sums to zero as complex numbers;
summing its multiplicities instead gives $2d=kd+p\sum_A n_A$.
Thus $p\mid d$ for $k=0,1$. The same identities prove sufficiency.

For sharpness take $d=mp$. On the standard torus one can prescribe all
character coordinates independently. For $k=0$ prescribe $m$ coordinates
equal to $1$ and $2m$ coordinates from each nontrivial inverse pair in
$H$, choosing either member in each coordinate. This uses $mp$ entries
and gives $c(x)=2m$ throughout $H$. For $k=1$ prescribe
$m(p+1)/2$ coordinates equal to $1$ and $m$ from each nontrivial inverse
pair in $H$. This gives $c(1)=d+m$, $c(x)=m$ for $x\in H\setminus\{1\}$.
\end{proof}

\begin{corollary}\label{c:dim}
For the rectangular torus with $d$ side lengths $L_i$ that are powers of
$p$, suppose $p\nmid d$. With $V=\prod_i L_i$ and $g=\gcd_i L_i$, the
periodic subspace has dimension
\[
 2(d-1)V+2+2^d(g-1)
\]
and the walk restricted to this subspace has exact order $2g$.
\end{corollary}
\begin{proof}
For each $u\in\mu_g\setminus\{1\}$, precisely $2^d$ tuples in
$\{u,u^{-1}\}^d$ contribute $u$ once. The trivial tuple contributes $1$
twice. Add the constant factors of \eqref{eq:fourier}. The eigenvalues
include $-1$ and all $g$th roots, giving order $2g$ since $g$ is odd.
\end{proof}

\section{Growth of the additional multiplicities}
For the standard torus $G=(\Z/N\Z)^d$, let $M_k(N)$ be the number of
characters with $T=k$ for $k=0,1$. This is also the multiplicity of
\emph{each} of the corresponding two roots, since they are distinct.

\begin{theorem}\label{t:count}
Let $d=mp$ for an odd prime $p$. Put $B=(N/p-1)/2$, $v=m\bmod 2$
with $v\in\{0,1\}$, and $s=\lfloor m/2\rfloor$. For $N=p^h$, $h\ge1$,
$M_0(N)$ and $M_1(N)$ are polynomials in $B$ of degrees $m$ and $s$.
Their leading terms in the binomial polynomial basis are respectively
\[
 2^d d!\binom Bm,\qquad
 \frac{2^{ps+(p-1)v/2}d!}{((d+v)/2)!}\binom Bs .
\]
In particular $M_0(N)\sim 2^d d!N^m/\bigl(m!(2p)^m\bigr)$.
\end{theorem}
\begin{proof}
Outside $H$ there are $B$ inverse-paired coset pairs, each of which
contains $p$ inverse pairs of elements. For $T=k$ each of these $p$
pairs must contribute an identical number $t_j$ of coordinates, because
each coordinate on an inverse pair contributes one to each of its two counts. In the
central coset, each of the $(p-1)/2$ nonidentity inverse pairs has
$r$ coordinates, and the identity has $a_0=(r+kd)/2$. Thus
\[
 d=a_0+(p-1)r/2+p\sum_{j=1}^B t_j .
\]
For $k=0$ this is $m=r/2+\sum t_j$, and $r$ is even. For $k=1$
it is $m=r+2\sum t_j$. These conditions are sufficient by
Theorem~\ref{t:main}. No aligned tuple has either trace, as $N$ is odd
and the roots belonging to these traces have even order.

For fixed $r$ and an ordered list of the positive $t_j$ on $b$ selected
coset pairs (in any fixed order), the number of tuples is
\begin{equation}\label{eq:weights}
 \binom Bb\,
 \frac{d!\,2^{d-a_0}}
 {a_0!\,(r!)^{(p-1)/2}\,\prod_{j=1}^{b}(t_j!)^{p}}.
\end{equation}
Here the multinomial factor assigns coordinates and each nonidentity
coordinate allows two choices of sign. Summing \eqref{eq:weights} over all
admissible $r,b$ and ordered positive lists is an exact formula. It
produces a polynomial since $b\le m$ for $k=0$ and $b\le s$ for $k=1$.
Equality in the first bound requires $r=0$ and all positive $t_j=1$.
Equality in the second requires $r=v$ and all positive $t_j=1$:
if $m$ is odd, $r$ must be odd. Substitution gives both claimed terms.
\end{proof}

For example, take $p=3$, $d=6$ and $B=(N/3-1)/2$. The complete formulas are
\[
 M_0(N)=240+17280B+46080\binom B2,\qquad M_1(N)=60+960B .
\]
At $N=9$ there are $17520$ copies of each primitive fourth root and
$1020$ of each primitive sixth root. The number of aligned nontrivial
quadratic modes grows only linearly in $N$ for fixed $d$. This family
therefore exhibits an exceptional quadratic contribution of order $N^2$
in a tower whose constant modes have order $N^6$.
The accompanying integer-arithmetic script checks these multiplicities
and tests the classification against cyclotomic reduction.

\section{Two pairs of steps at arbitrary exponent}
For a two-generated finite abelian group, the parity restriction on the
exponent can be removed, and the answer takes a different form.
\begin{theorem}\label{t:two}
For $d=2$, write $X=X_1,Y=X_2$. The complete periodic quadratic blocks,
with overlaps counted once, are as follows.
\begin{center}\small
\begin{tabular}{ll}\toprule
Condition & Quadratic eigenvalues\\\midrule
$Y\in\{X,X^{-1}\}$ & $X,X^{-1}$\\
$Y\in\{-X,-X^{-1}\}$ & $i,-i$\\
One coordinate equals $1$, the other is $i$ or $-i$ & primitive sixth roots\\
One coordinate equals $-1$, the other is $i$ or $-i$ & primitive third roots\\
\bottomrule\end{tabular}
\end{center}
\end{theorem}
\begin{proof}
Only primes over $2$ enter the integrality test for $(C(X)+C(Y))/2$.
Modulo the ideal $(2)$ the numerator equals
$C(X)-C(Y)=(X-Y)(X-Y^{-1})/X$. Normalize $v(2)=1$ at any prime above
$2$ in a cyclotomic field containing $X,Y$. For a root of unity $u$,
$v(1-u)=+\infty$ when $u=1$, is $2^{1-a}$ when its order is
$2^a>1$, and is zero when its order has a nontrivial odd part. Indeed
odd-order roots retain their order in residue characteristic two (apply
the geometric sum to any nontrivial odd-order root reducing to one);
$2$-power roots reduce to one. For a primitive $2^a$th root $u$, all
$1-u^j$ with $j$ odd have equal valuation by the geometric sum, and
their product is $\Phi_{2^a}(1)=2$ over the odd classes modulo $2^a$.
Consequently the test is $v(1-X/Y)+v(1-XY)\ge1$. It holds if either
argument $X/Y,XY$ is $1$ or $-1$. Otherwise it holds exactly when both
are primitive fourth roots, as all the remaining summands are at most
$1/2$. In that case one of $X,Y$ equals $\pm1$ and the other $\pm i$:
their squares are respectively the product and the quotient of those two
primitive fourth roots and have opposite signs. Substitution in
\eqref{eq:fourier} gives the table.
\end{proof}
For example on a $4\times4$ torus all sixteen character tuples qualify.
The four tuples with $X,Y\in\{i,-i\}$ also obey the second
condition and contribute their two roots just once. The walk is periodic
of exact order $12$. For odd rectangular tori only the first row remains,
giving the repository result for arbitrary odd side lengths.

\begin{thebibliography}{9}
\bibitem{repo} V. Reshetnikov, \emph{ProveIt}, Analysis/GeneralizedRiemannHypothesis,
Article, section on graph towers, subsection on exact spectral tests for
Grover walks, proposition on an odd rectangular grid.
\href{https://github.com/VladimirReshetnikov/ProveIt/blob/4ccd0fdeda0b1cc316d2dfdccfdbcfea9d4b77fb/Analysis/GeneralizedRiemannHypothesis/Article/sections/gra.tex}{Pinned source, commit 4ccd0fd}.
\bibitem{ahs} J. Akahori, T. Hayashi and R. Suzuki, Iwasawa-type spectral
resultant growth laws for Grover walks on graph towers, \emph{Ann. Math.
Qu\'ebec} (2026), \href{https://doi.org/10.1007/s40316-026-00290-1}{doi:10.1007/s40316-026-00290-1}.
\bibitem{speyer} D. E. Speyer, answer on the rational-angle cosine product
equation, MathOverflow (2023),
\href{https://mathoverflow.net/questions/458827}{question 458827}.
\bibitem{washington} L. C. Washington, \emph{Introduction to Cyclotomic
Fields}, 2nd ed., Springer GTM 83 (1997).
\end{thebibliography}
\end{document}
